\chapter{Chercheur quantitatif}\label{ch:in:quant-researcher}

Un chercheur dont le signal est négocié chaque jour et dont le ratio de Sharpe atteint 2 peut établir, en un an, que le résultat n'est pas
dû au hasard. Un chercheur dont le signal est rééquilibré chaque mois et dont le ratio de Sharpe atteint un demi a besoin de quinze ans et demi de
résultats en direct pour établir la même chose. Les deux portent le même titre et peuvent travailler dans le même immeuble, mais l'un apprend du
marché chaque mois, l'autre peut-être deux fois dans une carrière. Ce chapitre décrit le \omterm{def:in:quant-researcher:def}{chercheur quantitatif} chez chaque
type d'employeur, mesure la vitesse à laquelle le travail du poste peut être jugé, et suit ce dont un chercheur est propriétaire, ce qu'il livre et
ce dont il provient~; les éléments sur la rémunération du chapitre~14 s'y ajoutent.

\begin{center}\footnotesize
\begin{tabular}{@{}p{3.2cm}p{9.4cm}@{}}
\toprule
\multicolumn{2}{@{}l}{\textbf{Fiche de poste~: \omterm{def:in:quant-researcher:def}{chercheur quantitatif}}}\\
\midrule
employeurs & fonds systématiques, teneurs de marché, plateformes, banques, sociétés de gestion\\
horizon des décisions & de quelques secondes à plusieurs mois\\
P\&L & le partage, par le mérite attribué aux signaux\\
rattaché à & le responsable de la recherche ou un gérant de portefeuille\\
codes de profession dans les dépôts & 13-2099.01, 15-2041, 15-2031~; intitulés contenant «~quantitative research~», «~quantitative analyst~»\\
où la série l'enseigne & livre~7, chapitres~1 et~15~; livre~16, chapitre~9\\
\bottomrule
\end{tabular}
\end{center}

\section{Trois métiers de recherche~: haute fréquence, moyenne fréquence, société de gestion}

\begin{definition}[Chercheur quantitatif]\label{def:in:quant-researcher:def}
Un \emph{chercheur quantitatif}\index{chercheur quantitatif} est une personne dont le travail consiste à trouver, tester et maintenir
des régularités statistiques qu'une société peut négocier ou utiliser pour gérer le risque, et à les transformer en spécifications, modèles ou
code que les systèmes de négociation de la société exécutent.
\end{definition}

Un teneur de marché décrit le rôle dans ses propres termes~: les \omterm{def:in:quant-researcher:def}{chercheurs quantitatifs} «~se concentrent sur la résolution de problèmes complexes
par une recherche approfondie et indépendante. Leur travail façonne les stratégies de négociation, améliore l'exécution et optimise la gestion des positions.
Certains projets naissent d'une idée rapide, d'autres demandent des investigations plus longues et des itérations.~» Trois
cadres donnent à ces mots des contenus différents.
\begin{itemize}
\item \textbf{Haute fréquence} (chapitre~2)~: des signaux sur des secondes à des minutes, tirés des données du carnet d'ordres~; le travail est
proche du système de négociation, et un changement se juge en quelques jours sur des milliers de transactions indépendantes (livre~11).
\item \textbf{Moyenne fréquence} (chapitres~4 et~6)~: des signaux sur des jours à des semaines, tirés des prix, des données fondamentales et des
données alternatives~; un signal rejoint un portefeuille avec d'autres et se juge sur plusieurs mois (livres~7 et~8).
\item \textbf{Gestion d'actifs} (chapitre~8)~: des modèles de facteurs et d'allocation sur des mois à des années, où l'essentiel des
preuves vient de l'historique, et les résultats en direct arrivent lentement (livre~8).
\end{itemize}
Les analystes quantitatifs d'une banque travaillent surtout sur les modèles de prix et de risque plutôt que sur l'alpha~; le chapitre~18 les traite.

\section{Horizon, données et vitesse du retour d'information}

\begin{definition}[Vitesse du retour d'information]\label{def:in:quant-researcher:feedback}
La \emph{vitesse du retour d'information}\index{vitesse du retour d'information} d'un poste de recherche est le rythme auquel son travail produit des preuves
capables de distinguer le talent de la chance~: le temps calendaire nécessaire pour qu'un résultat de taille donnée devienne
statistiquement distinguable de zéro.
\end{definition}

\begin{method}[Temps jusqu'à la significativité]\label{met:in:quant-researcher:time}
Soit $S$ le ratio de Sharpe annuel vrai d'une stratégie, rééquilibrée $m$ fois par an, avec des rendements indépendants. Le ratio de Sharpe estimé
par période a une erreur type $\sqrt{(1+s^2/2)/n}$ après $n$ périodes, où $s=S/\sqrt m$ (Lo, 2002). Le
nombre de périodes nécessaires pour que l'estimation soit à $z$ erreurs types de zéro est
\[
n=\frac{z^2\,(1+s^2/2)}{s^2},
\]
et le temps calendaire vaut $n/m$ années. Avec $z=1.96$, $n$ est proche de $3.84\,m/S^2$ périodes, et le temps calendaire proche de
$3.84/S^2$ années.
\end{method}

La formule dit que le temps calendaire dépend du ratio de Sharpe, non de l'horizon~: un ratio de Sharpe de 2 demande environ un an, de 1 environ
quatre ans, de un demi environ quinze. L'horizon intervient par le ratio de Sharpe qu'un signal peut atteindre.
Par la loi fondamentale de la gestion active (Grinold, 1989~; livre~7, chapitre~15), un signal de talent $\mathrm{IC}$ par
pari indépendant et $N$ paris indépendants par période rapporte $S\approx \mathrm{IC}\sqrt{N m}$~: des horizons plus courts donnent plus de paris
par an, donc des ratios de Sharpe plus élevés, pour un même talent par pari.

\begin{example}[Le même talent à six horizons]\label{ex:in:quant-researcher:law}
Prenons un talent illustratif de 0.02 par pari, sur 50 paris indépendants à chaque rééquilibrage. Le temps jusqu'à la significativité est alors d'environ
194 périodes de rééquilibrage à chaque horizon, car $n\approx z^2/(\mathrm{IC}^2N)$ ne dépend pas de $m$~: 0.7 jour à un
horizon d'une minute, 43 jours à une heure, 0.77 an à un jour, 3.73 ans à une semaine, 16.2 ans à un mois et 48.5
ans à un trimestre (\cref{fig:in:quant-researcher:feedback}).
\end{example}

\begin{omfigure}
\begin{tikzpicture}
\begin{axis}[width=10.5cm,height=6cm,xmode=log,ymode=log,xmin=0.001,xmax=100,ymin=0.001,ymax=100,
  xlabel={\small horizon de rééquilibrage, jours de négociation (échelle log)},ylabel={\small années jusqu'à la significativité (échelle log)},
  tick label style={font=\footnotesize},grid=major,grid style={gray!20},table/col sep=comma,
  legend style={font=\scriptsize,at={(0.5,-0.3)},anchor=north,legend cell align=left}]
\addplot[omDef,thick,mark=*] table[x=days,y=years]{figdata/industry/17-quant-researcher/feedback.csv};
\addplot[only marks,mark=square*,mark size=2.5pt,omThm] table[x=days,y=years]{figdata/industry/17-quant-researcher/named.csv};
\legend{même talent par pari (IC 0.02~; 50 paris par période),Sharpe 2 quotidien~; 1 hebdomadaire~; 0.5 mensuel}
\end{axis}
\end{tikzpicture}

\omcaption{Temps calendaire avant que le ratio de Sharpe d'une stratégie ne soit distinguable de zéro à 95~\%, en fonction de son horizon de rééquilibrage
(une minute à un trimestre). La courbe maintient le talent par pari fixe~; les carrés sont les trois stratégies nommées du chapitre
stratégies. Le talent et la largeur sont illustratifs. Données~: \texttt{in\_researcher.law\_curve} et \texttt{named}.}\label{fig:in:quant-researcher:feedback}
\end{omfigure}

La \omterm{def:in:quant-researcher:feedback}{vitesse du retour d'information} façonne le métier. Un chercheur en haute fréquence peut essayer beaucoup d'idées par an et apprendre lesquelles fonctionnent~; un
chercheur en moyenne fréquence doit s'appuyer davantage sur l'historique, sur la discipline hors échantillon et sur les contrôles de recherche du livre~7,
car le marché ne tranchera pas à temps~; un chercheur en gestion d'actifs verra rarement un résultat en direct qui prouve un modèle
avant que celui-ci n'ait été utilisé pendant des années. Elle façonne aussi les carrières~: le chercheur dont les résultats sont jugés vite peut construire un
historique vite (chapitre~28), et la rémunération peut suivre les résultats de plus près (chapitre~13).

\section{Propriété du P\&L et du mérite}

Un chercheur ne possède guère de book. Le P\&L appartient à un gérant de portefeuille, à un desk ou à la société, et le
chercheur est rémunéré pour y contribuer. La mesure de cette contribution relève de l'attribution du mérite de la société (livre~16,
chapitre~9)~: les signaux d'un portefeuille combiné rapportent ensemble, et répartir leur profit commun entre leurs auteurs est un
choix, non un fait. Les trois cadres diffèrent.
\begin{itemize}
\item Dans une plateforme (chapitre~5), le chercheur travaille pour un gérant de portefeuille dont le book est jugé sur son propre résultat~; le
mérite du chercheur relève de la décision du gérant.
\item Dans un fonds systématique (chapitre~4), un signal rejoint un portefeuille de recherche géré de façon centralisée~; le mérite suit la
méthode d'attribution de la société, et la société possède le signal.
\item Chez un teneur de marché (chapitres~2--3), la recherche améliore la cotation et la couverture qu'utilisent de nombreux traders~; le mérite est partagé par
conception, et la rémunération suit le résultat du desk ou de la société.
\end{itemize}
Un candidat doit demander comment le mérite est mesuré et qui en décide~; les structures du chapitre~13 reposent sur cette réponse.

\section{Ce que livre un chercheur}

Le produit d'un chercheur n'est pas un article. C'est une modification que les systèmes de la société exécutent~: un signal avec sa spécification, les
données dont il a besoin et la manière de le calculer~; les tests et le journal de recherche qui le justifient (livre~7, chapitre~1)~; du code qui
respecte les standards de la société, ou une spécification que les développeurs implémentent (chapitre~19)~; et un plan de suivi qui indique comment
la société saura qu'il a cessé de fonctionner. Selon les mots du teneur de marché, les chercheurs «~conçoivent et optimisent des algorithmes de trading~».
Les revues de recherche décident de ce qui est livré~; la plupart de ce qu'un chercheur essaie ne l'est pas.

\section{D'où viennent les chercheurs}

La classification professionnelle que les dépôts utilisent pour l'essentiel du poste, 13-2099.01 «~Financial Quantitative
Analysts~», figure dans la plus haute bande de préparation d'O*NET~: «~Most of these occupations require graduate
school. For example, they may require a master's degree, and some require a Ph.D.~» En exercice 2025, 1\,018 des
candidatures classées dans la famille de postes portaient ce code, 156 statisticiens (15-2041) et 52 analystes en recherche
opérationnelle (15-2031)~; le poste fait appel aux mathématiques, à la statistique, à la physique, à l'informatique et à l'économie, et le chapitre~29
décrit les filières.

\begin{dated}{2025-09}{Ce que gagnent les chercheurs quantitatifs~: les données publiques}\label{dat:in:quant-researcher:pay}
Demandes de certification de conditions de travail américaines, exercice 2025, famille de postes «~quantitative researcher or analyst~» (intitulés, ou SOC
13-2099.01 lorsque l'intitulé est générique)~: \omterm{def:in:how-pay-works:base}{salaire de base} médian proposé de 220\,000~\$ chez les fonds systématiques (120 candidatures, 5 employeurs),
190\,000~\$ chez les plateformes (45, 3), 175\,000~\$ chez les teneurs de marché (131, 9), 158\,100~\$ chez les banques (825, 9) et 107\,100~\$ chez les
bourses (35, 4). Par \omterm{def:in:pay-levels-by-role-firm-type-and-seniority:soc}{niveau de salaire} chez les banques~: 88\,300~\$ (I), 145\,300~\$ (II), 179\,335~\$ (III), 200\,000~\$ (IV)~; chez les
teneurs de marché, 150\,000~\$, 200\,000~\$, 150\,000~\$, 225\,000~\$. Enquête professionnelle sur les spécialistes financiers du
secteur des valeurs mobilières, mai 2025~: médiane de 100\,190~\$, entre le 10\up{e} et le 90\up{e}~centile, de 55\,370~\$ à 222\,100~\$. \omterm{def:in:how-pay-works:base}{Salaire de base} uniquement.
\end{dated}

\begin{omfigure}
\begin{tikzpicture}
\begin{axis}[width=10.5cm,height=5.2cm,xmin=50,xmax=300,ymin=0.4,ymax=5.6,ytick=data,
  yticklabels={{bourse},{banque},{teneur de marché},{plateforme},{fonds systématique}},yticklabel style={font=\footnotesize},
  xticklabel style={font=\footnotesize},xlabel={\small salaire de base proposé, milliers de \$},grid=major,grid style={gray!20},
  table/col sep=comma]
\addplot[draw=none,mark=none,error bars/.cd,x dir=both,x explicit,error mark=none,error bar style={omDef,line width=1pt}]
  table[x=p50,y=pos,x error plus expr=\thisrow{p90}-\thisrow{p50},x error minus expr=\thisrow{p50}-\thisrow{p10}]
  {figdata/industry/17-quant-researcher/pay_kind.csv};
\addplot[draw=none,mark=none,error bars/.cd,x dir=both,x explicit,error mark=none,error bar style={omDef!40,line width=7pt}]
  table[x=p50,y=pos,x error plus expr=\thisrow{p75}-\thisrow{p50},x error minus expr=\thisrow{p50}-\thisrow{p25}]
  {figdata/industry/17-quant-researcher/pay_kind.csv};
\addplot[only marks,mark=|,mark size=5pt,line width=1.5pt,omThm] table[x=p50,y=pos]{figdata/industry/17-quant-researcher/pay_kind.csv};
\end{axis}
\end{tikzpicture}

\omcaption{\omterm{def:in:how-pay-works:base}{Salaires de base} proposés aux \omterm{def:in:quant-researcher:def}{chercheurs quantitatifs} dans les demandes de certification de l'exercice 2025, par type d'employeur~: du 10\up{e} au
90\up{e}~centile (fin), intervalle interquartile (épais), médiane (repère). Données~: \texttt{data/industry/lca\_ranges.csv}, par l'intermédiaire de
\texttt{in\_researcher.pay}.}\label{fig:in:quant-researcher:pay}
\end{omfigure}

L'ordre des employeurs est celui du chapitre~14 pour tous les postes. Les dépôts des banques montrent l'échelle de séniorité la plus nette
(\cref{fig:in:quant-researcher:levels})~; ceux des teneurs de marché non, pour la raison donnée au chapitre~14~: les offres très supérieures au
\omterm{def:in:pay-levels-by-role-firm-type-and-seniority:lca}{salaire en vigueur} ne sont pas ordonnées par le niveau que l'employeur déclare.

\begin{omfigure}
\begin{tikzpicture}
\begin{axis}[width=10.5cm,height=5cm,xmin=0.8,xmax=4.2,ymin=60,ymax=250,xtick={1,2,3,4},xticklabels={I,II,III,IV},
  xlabel={\small niveau de salaire},ylabel={\small salaire de base médian proposé, milliers de \$},tick label style={font=\footnotesize},
  grid=major,grid style={gray!20},table/col sep=comma,unbounded coords=jump,
  legend style={legend cell align=left,font=\scriptsize,at={(0.5,-0.3)},anchor=north,/tikz/every even column/.append style={column sep=8pt}},
  legend columns=3]
\addplot[omThm,thick,mark=square*] table[x=level,y=market_maker]{figdata/industry/17-quant-researcher/levels.csv};
\addplot[omDef,thick,mark=*] table[x=level,y=bank]{figdata/industry/17-quant-researcher/levels.csv};
\addplot[black!70,thick,mark=triangle*] table[x=level,y=systematic_fund]{figdata/industry/17-quant-researcher/levels.csv};
\legend{teneur de marché,banque,fonds systématique}
\end{axis}
\end{tikzpicture}

\omcaption{\omterm{def:in:how-pay-works:base}{Salaire de base} médian proposé aux \omterm{def:in:quant-researcher:def}{chercheurs quantitatifs} selon le \omterm{def:in:pay-levels-by-role-firm-type-and-seniority:soc}{niveau de salaire}, exercice 2025. Les cellules supprimées sont omises.
Données~: comme la \cref{fig:in:quant-researcher:pay}.}\label{fig:in:quant-researcher:levels}
\end{omfigure}

\section{Tutoriel~: combien de temps avant de savoir~?}\label{tut:in:quant-researcher:tut}

\begin{tutorial}
\textbf{Objectif.} Chiffrer la \omterm{def:in:quant-researcher:feedback}{vitesse du retour d'information} du poste et la placer à côté des données de salaire.
\textbf{État final~:} \cref{fig:in:quant-researcher:feedback,fig:in:quant-researcher:pay,fig:in:quant-researcher:levels}.
\begin{tutsteps}
\item \textbf{Salaire.} \texttt{firm.roles.lca\_cells(rows, role)} avec le poste \texttt{quant researcher}, à tous les niveaux et par
niveau~; \texttt{data/industry/lca\_titles.csv} donne les codes professionnels derrière le poste.
\item \textbf{Temps jusqu'à la significativité.} \texttt{periods\_to\_significance(S, m)} et \texttt{years\_to\_significance}
(\cref{lst:in:quant-researcher:time}).
\omcode{code/firm/roles/firm_roles.py}{124}{132}{Nombre de périodes et d'années de données avant qu'un ratio de Sharpe se distingue de
zéro, avec l'erreur type de Lo.}\label{lst:in:quant-researcher:time}
\item \textbf{La loi.} \texttt{fundamental\_law\_sr(IC, N, m)} transforme la compétence et la largeur en ratio de Sharpe pour chaque
horizon.
\end{tutsteps}
\end{tutorial}

Pour les trois stratégies du chapitre, un ratio de Sharpe de 2 à horizon quotidien demande 0.97 année (244 jours de transaction)~; 1 à horizon
hebdomadaire, 3.88 années (202 semaines)~; la moitié à horizon mensuel, 15.53 années (186 mois).

\textbf{Pour aller plus loin.} Supposer les rendements autocorrélés et utiliser l'erreur type ajustée de Lo~; compter les essais qu'un
chercheur a menés et déflater le ratio de Sharpe en conséquence (livre~7).

\section{Construction~: la fiche du chercheur et la vitesse du retour d'information}\label{bld:in:quant-researcher:roles}

\begin{build}
\bfield{Rôle} Ajouter le \omterm{def:in:quant-researcher:def}{chercheur quantitatif} au registre des postes, ainsi que les fonctions qui mesurent la vitesse à laquelle la recherche
est jugée.
\bfield{Interface} \texttt{firm.roles}~: la fiche \texttt{quant researcher}~;
\texttt{periods\_to\_significance(sr, periods\_per\_year, z)}~; \texttt{years\_to\_significance}~;
\texttt{fundamental\_law\_sr(ic, breadth, periods\_per\_year)}.
\bfield{Règles} Les rendements sont IID, sauf si l'appelant les ajuste~; le ratio de Sharpe est annualisé~; la compétence et la largeur sont celles de l'appelant,
étiquetées à titre illustratif.
\bfield{Tests d'acceptation} \texttt{code/firm/roles/tests/}~: le temps jusqu'à la significativité vaut $3.84/S^2$ années au premier ordre et
exactement la valeur de la formule~; il décroît avec le ratio de Sharpe~; sous la loi fondamentale, le nombre de périodes est le même à
chaque horizon.
\bfield{Pour aller plus loin} Autocorrélation et rendements non gaussiens~; déflation pour tests multiples~; une version bayésienne qui met à jour
un a priori sur la compétence.
\end{build}

\begin{omsources}
\item Lo, A.~W. (2002), The statistics of Sharpe ratios, \emph{Financial Analysts Journal} 58(4), 36--52.
\item Grinold, R.~C. (1989), The fundamental law of active management, \emph{Journal of Portfolio Management} 15(3), 30--37.
\item Optiver, « Breaking down the trading industry »~; O*NET, 13-2099.01.
\item Tableaux du chapitre~14 issus des fichiers LCA du Department of Labor et de l'enquête professionnelle du BLS.
\end{omsources}

\section{Exercices}

\begin{exercise}[$\star$]\label{exo:in:quant-researcher:1}
Combien d'années de rendements quotidiens faut-il pour qu'un vrai ratio de Sharpe de 1.5 se distingue de zéro à 95~\%~?
\end{exercise}

\begin{exercise}[$\star$]\label{exo:in:quant-researcher:2}
Avec la loi fondamentale, une compétence de 0.02 et 50 paris par période, quel est le ratio de Sharpe annuel à un horizon hebdomadaire~?
\end{exercise}

\begin{exercise}[$\star$]\label{exo:in:quant-researcher:3}
D'après l'encadré daté, quel type d'employeur offre aux \omterm{def:in:quant-researcher:def}{chercheurs quantitatifs} le \omterm{def:in:how-pay-works:base}{salaire de base} médian le plus élevé, et de
combien dépasse-t-il celui des banques~?
\end{exercise}

\begin{exercise}[$\star\star$]\label{exo:in:quant-researcher:4}
Montrer que, sous la loi fondamentale, le nombre de périodes nécessaires pour atteindre la significativité ne dépend pas de l'horizon.
\end{exercise}

\begin{exercise}[$\star\star$]\label{exo:in:quant-researcher:5}
Le signal mensuel d'un chercheur affiche un ratio de Sharpe de 0.8 sur trois ans de trading en production. Est-ce une preuve de compétence~?
\end{exercise}

\begin{exercise}[$\star\star$]\label{exo:in:quant-researcher:6}
Pourquoi la reconnaissance accordée à un chercheur en plateforme dépend-elle davantage du jugement d'une seule personne que celle d'un chercheur en fonds systématique~?
\end{exercise}

\begin{exercise}[$\star\star\star$]\label{exo:in:quant-researcher:7}
\emph{Programmation.} Des rendements à autocorrélation d'ordre un $\rho$ gonflent la variance du ratio de Sharpe annualisé
d'environ $(1+\rho)/(1-\rho)$. Recalculer le temps nécessaire à la significativité de la stratégie hebdomadaire avec $\rho=0.2$.
\end{exercise}

\begin{exercise}[$\star\star\star$]\label{exo:in:quant-researcher:8}
\emph{Trouver la faille.} «~Notre stratégie mensuelle affiche un ratio de Sharpe de 1.2 en backtest sur 20 ans, elle sera donc validée
en production dans les trois ans.~»
\end{exercise}

\section{Problème~: Combien de temps avant de savoir~?}

\begin{problem}[{Problème du week-end --- combien de temps avant de savoir~?}]\label{pb:in:quant-researcher:1}
Une diplômée ayant reçu des offres en recherche à haute fréquence et dans un fonds systématique à rééquilibrage mensuel
se demande en combien de temps chaque poste lui dirait si elle y est douée.

\medskip\noindent\textbf{Partie I --- Le rôle.}
\begin{enumerate}
\item Définir un \omterm{def:in:quant-researcher:def}{chercheur quantitatif}~; décrire les trois contextes de recherche.
\item Comment un teneur de marché décrit-il le rôle~?
\item En quoi le travail quantitatif d'une banque diffère-t-il~?
\item Définir la \omterm{def:in:quant-researcher:feedback}{vitesse du retour d'information}.
\item Énoncer la fiche de rôle.
\end{enumerate}

\medskip\noindent\textbf{Partie II --- Les statistiques.}
\begin{enumerate}[resume]
\item Énoncer l'erreur type du ratio de Sharpe de Lo sous des rendements IID.
\item Établir le nombre de périodes nécessaires à la significativité.
\item Pourquoi le temps calendaire dépend-il du ratio de Sharpe plutôt que de l'horizon~?
\item Énoncer la loi fondamentale et ce qu'elle implique pour l'horizon.
\item Donner les durées pour les six horizons de l'exemple.
\end{enumerate}

\medskip\noindent\textbf{Partie III --- Le poste.}
\begin{enumerate}[resume]
\item Qui porte le P\&L sur une plateforme, dans un fonds systématique et chez un teneur de marché~?
\item Que livre un chercheur~?
\item D'où viennent les chercheurs, et quelle préparation le métier exige-t-il~?
\item Que montrent les données de rémunération par employeur et par niveau~?
\item Pourquoi les dépôts des teneurs de marché ne montrent-ils pas d'échelle de séniorité~?
\end{enumerate}

\medskip\noindent\textbf{Partie IV --- Le verdict.}
\begin{enumerate}[resume]
\item Énoncer le \emph{résultat nommé}~: le temps nécessaire pour rejeter un ratio de Sharpe nul, pour un ratio de Sharpe de 2 à horizon d'un jour,
de 1 à une semaine et de 0.5 à un mois.
\item Que ferait l'autocorrélation à ces durées~?
\item Que peut faire un chercheur dans un poste à retour d'information lent pour apprendre plus vite~?
\item En quoi cette différence doit-elle influencer la manière dont elle juge ses deux premières années dans chaque poste~?
\item En deux phrases, lui répondre.
\end{enumerate}
\end{problem}

\section{Questions d'entretien}

\begin{interviewq}[$\star$ \iqroles{chercheur}]\label{iq:in:quant-researcher:1}
Un estimé quotidien du ratio de Sharpe d'une stratégie vaut 0.1 sur 400 jours. Quelle est sa valeur annualisée, et à peu près son erreur type~?
\end{interviewq}

\begin{interviewq}[$\star$ \iqroles{chercheur, trader}]\label{iq:in:quant-researcher:2}
Pourquoi les stratégies à haute fréquence affichent-elles souvent des ratios de Sharpe élevés et une faible capacité~?
\end{interviewq}

\begin{interviewq}[$\star\star$ \iqroles{chercheur}]\label{iq:in:quant-researcher:3}
Vous avez testé 200 signaux et le meilleur affiche une t-statistique de 3.5 en backtest. Que concluez-vous~?
\end{interviewq}

\begin{interviewq}[$\star\star$ \iqroles{chercheur, ingénieur ML}]\label{iq:in:quant-researcher:4}
Comment décideriez-vous qu'un signal en production a cessé de fonctionner~?
\end{interviewq}

\begin{interviewq}[$\star\star$ \iqroles{chercheur, développeur}]\label{iq:in:quant-researcher:5}
Que doit contenir la spécification d'un signal pour qu'un développeur puisse l'implémenter sans vous~?
\end{interviewq}

\begin{interviewq}[$\star\star\star$ \iqroles{chercheur}]\label{iq:in:quant-researcher:6}
Deux signaux partagent 60~\% de leurs positions. Comment répartiriez-vous le mérite de leur profit commun~?
\end{interviewq}
